\documentclass[11pt,a4paper]{article}

% === Typography & Language ===
\usepackage[utf8]{inputenc}
\usepackage[T1]{fontenc}
\usepackage[english]{babel}
\usepackage{microtype} % Implements micro-typographic extensions for better spacing
\usepackage{charter}   % Professional, highly readable serif font for engineering documents

% === Page Layout ===
\usepackage[margin=1in, bindingoffset=0mm]{geometry}
\usepackage{fancyhdr}  % Custom headers and footers
\usepackage{parskip}   % Modern paragraph spacing (adds space instead of indenting)
\usepackage{float}
\usepackage{needspace}
\usepackage{lastpage}
\usepackage{lineno}    % Line numbering for review drafts

% === Colors & Visual Styling ===
\usepackage{xcolor}
\definecolor{NavyBlue}{HTML}{1A365D}    % Primary Accent: Deep Navy
\definecolor{SlateGrey}{HTML}{4A5568}   % Secondary Accent: Cool Grey
\definecolor{LightCyan}{HTML}{EDF2F7}   % Light background tint for callouts
\definecolor{BorderGrey}{HTML}{CBD5E0}  % Subtle border color

% === Tables & Lists ===
\usepackage{booktabs}   % Publication-quality tables
\usepackage{tabularx}   % Tables with adjustable-width columns
\usepackage{enumitem}   % Customized lists
\usepackage{amssymb}    % Required for checkbox symbols ($\square$)

% === Acronyms ===
\usepackage{acronym}

% === Custom Environments ===
\newenvironment{deflist}{%
  \begin{description}[
    font=\normalfont\bfseries,
    leftmargin=1.5em,
    labelindent=0em,
    itemsep=4pt,
    parsep=0pt
  ]
}{%
  \end{description}
}

\newif\ifdraft
\drafttrue  % <-- change to \draftfalse to publish

\ifdraft
  \usepackage{draftwatermark}
  \SetWatermarkText{DRAFT}
  \SetWatermarkScale{1}
  \SetWatermarkColor[gray]{0.85}
\fi

\usepackage[skins]{tcolorbox}
\newtcolorbox{principle}[1]{%
  enhanced,
  frame hidden,
  borderline west={3pt}{0pt}{NavyBlue},
  colback=LightCyan,
  sharp corners,
  left=10pt,
  right=8pt,
  top=7pt,
  bottom=7pt,
  before upper={\textbf{\color{NavyBlue}#1}\par\smallskip}%
}

% === Section Headings Customization ===
\usepackage{titlesec}
\titleformat{\section}{\large\bfseries}{\thesection}{1em}{}[{\color{BorderGrey}\titlerule}]
\titleformat{\subsection}{\normalsize\bfseries}{\thesubsection}{1em}{}

% === Hyperlinks & Cross-Referencing ===
\usepackage[hidelinks]{hyperref}

% === Bibliography ===
\usepackage[style=ieee,backend=biber]{biblatex}
\addbibresource{references.bib}

% === Header and Footer Setup ===
\pagestyle{fancy}
\fancyhf{}
\renewcommand{\headrulewidth}{0.4pt}
\renewcommand{\footrulewidth}{0.4pt}
\fancyhead[L]{\small Data-Centric Engineering AI Ethical Framework}
%\fancyhead[R]{\small Framework White Paper}
%\ifdraft\fancyfoot[L]{\small\textcolor{red}{DRAFT --- Not for Distribution}}\fi
\fancyfoot[C]{\small Page \thepage\ of \pageref{LastPage}}

% =========================================================================
% DOCUMENT METADATA
% =========================================================================
\title{
    \vspace{-1.5cm}
    \Huge\bfseries An AI Ethical Framework for Data-Centric Engineering \\
    %\Large\color{SlateGrey}Operationalizing Accountability, Traceability, and Boundaries in Safety-Critical Pipelines
}
\author{
    \textbf{Jacob Frusher} \\
    \small Education Intern, Department of Mechanical Engineering \\
    \small University of Southampton \\
    \small\href{mailto:jf7g23@soton.ac.uk}{\texttt{jf7g23@soton.ac.uk}}
}
\date{\small June 2026}

% =========================================================================
% DOCUMENT BODY
% =========================================================================
\begin{document}
\ifdraft\linenumbers\fi

% Acronym definitions --- printed list appears before the bibliography.
% Use \ac{key} in running text; first use auto-expands to full form.
\acrodef{AI}{Artificial Intelligence}
\acrodef{ASCE}{American Society of Civil Engineers}
\acrodef{ASME}{American Society of Mechanical Engineers}
\acrodef{CICD}[CI/CD]{Continuous Integration/Continuous Deployment}
\acrodef{DCE}{Data-Centric Engineering}
\acrodef{FMEA}{Failure Mode and Effects Analysis}
\acrodef{HITL}{Human-in-the-Loop}
\acrodef{IEC}{International Electrotechnical Commission}
\acrodef{IEEE}{Institute of Electrical and Electronics Engineers}
\acrodef{ISO}{International Organization for Standardization}
\acrodef{OOD}{Out-of-Distribution}

\maketitle

\begin{table}[H]
\centering
\small
\label{tab:doc_control}
\vspace{2mm}
\begin{tabularx}{\textwidth}{| l | X |}
\hline
\textbf{Document Reference} & UOS-MECH-DCE-2026-01 \\ \hline
\textbf{Current Version} & Version 1.0 \\ \hline
\textbf{Issue Date} & \today \\ \hline
\textbf{Originating Author} & Jacob Frusher \\ \hline
\end{tabularx}

\vspace{4mm}

\begin{tabularx}{\textwidth}{| l | l | X | l |}
\hline
\textbf{Rev.} & \textbf{Date} & \textbf{Summary of Changes} & \textbf{Author} \\ \hline
0.1 & June 14, 2026 & Initial drafting, scope definition, and foundational principle mapping. & J. Frusher \\ \hline
0.2 & June 18, 2026 & Integration of the operational lifecycle gates and tier architecture. & J. Frusher \\ \hline
0.3 & June 24, 2026 & Comprehensive revision to eliminate socio-technical transience; abstracted risk tiers around analytical tractability and verification. & J. Frusher \\ \hline
\end{tabularx}
\end{table}

\begin{abstract}
\noindent %\textbf{Abstract [Target Word Count: 150--200 words]} \\
The rapid integration of artificial intelligence into safety-critical engineering systems reveals a key regulatory gap: current ethical and governance frameworks for AI are rooted in software-centric paradigms that fall short when algorithmic failure carries direct physical, clinical, or infrastructural consequences. This article proposes an ethical framework geared towards engineers, developed within the particular setting of \ac{DCE}. The framework is structured around five core engineering principles, Verification, Traceability, Ultimate Liability, Boundary Clarity, and Stewardship, each conceived as a professional imperative binding the practising engineer rather than the algorithm. To ensure these principles are operationalized beyond philosophy, the framework introduces a proportional, three-tiered risk classification that scales compliance obligations to the stakes of each deployment context. These tiers underpin three sequential phases of implementation: an algorithmic risk gate at project inception, automated compliance checks embedded within \ac{CICD} pipelines, and a structurally decoupled runtime architecture mandating deterministically safe fallback and \ac{HITL} override for safety-critical systems. The framework further argues that effective governance requires a corresponding cultural shift in academic training, conditioning practitioners to treat model outputs as conditional probability distributions requiring validation rather than as verified truths. The overarching principle is that professional accountability cannot be delegated to a computational process.
\end{abstract}
\acresetall

\section{Introduction and Context}
\label{sec:intro}


\subsection{Background and Institutional Context}
This framework represents one of the key outputs of this research project, providing a formal ethical framework for the use of \ac{AI} in \ac{DCE}. It serves as an authoritative model providing a baseline for ensuring that deployed computational pipelines are verifiable, transparent, and traceable.. Taking into account the outstanding speed of current technological advancements, it is crucial to acknowledge that the proposed framework is designed to function as a living document, which can be iterated upon dynamically. It is critical to take preemptive measures due to the emergence of socio-technical risks associated with new types of engineering.

\subsection{The Engineering Gap in AI Ethics}
While established frameworks such as the \ac{IEEE} Global Initiative on Ethics of Autonomous and Intelligent Systems~\cite{ieee_ethically_aligned} and \ac{ISO}/\ac{IEC}~42001~\cite{iso42001} provide essential governance structures for data fairness and AI management systems, they are rooted in a software-centric paradigm. They do not address the cyber-physical realities codified in traditional engineering safety standards such as IEC~61508~\cite{iec61508} or ISO~13849~\cite{iso13849}. Furthermore, professional engineering codes of ethics, such as those of the \ac{ASME}~\cite{asme_ethics} and \ac{ASCE}~\cite{asce_ethics}, ground liability in licensed professional practice, a principle this framework extends to AI pipeline authorship. It is the aim of the proposed framework to bridge this vital regulatory gap.

\subsection{Terminology and Definitions}
\label{sec:terminology}
The following terms are used with specific technical meaning throughout this document. Precise definitions are necessary to prevent the semantic ambiguity that commonly undermines the practical application of ethical frameworks in engineering contexts.

\begin{deflist}

\item[Artificial Intelligence (AI).] As used in this framework, AI denotes any computational system that employs a statistical model, trained on empirical data, to perform prediction, classification, optimisation, or decision-support tasks. This definition explicitly excludes deterministic, rule-based algorithms, which carry conventional engineering liability without the additional epistemic risks addressed here.

\item[Data-Centric Engineering (DCE).] An engineering paradigm in which the data pipeline, its provenance, representativeness, and integrity, is treated as the primary engineering artefact, rather than the model architecture itself. DCE recognises that the performance and safety envelope of any AI system is fundamentally bounded by the quality and scope of its training data.

\item[Algorithmic Pipeline.] The end-to-end sequence of computational processes through which raw input data is transformed into a model output or engineering recommendation. This encompasses data ingestion, pre-processing, model inference, post-processing, and any downstream actuation or advisory stages.

\item[Out-of-Distribution (OOD) Data.] Input data whose statistical characteristics fall outside the distribution represented by the model's training dataset. OOD inputs invalidate the implicit assumptions of a trained model and constitute a primary and frequently underestimated failure mode in deployed AI systems~\cite{hendrycks2017}.

\item[Human-in-the-Loop (HITL).] A system architecture or operational protocol in which a qualified engineer is required to review, validate, or explicitly approve algorithmic outputs before consequential action is taken. HITL mechanisms are the primary instrument through which nondelegable professional liability is preserved at runtime.

\item[Continuous Integration / Continuous Deployment (CI/CD).] A software engineering practice in which code changes are automatically compiled, tested, and validated against a defined suite of checks prior to integration or deployment. Within this framework, the CI/CD pipeline serves as the enforcement mechanism for automated ethical and physical boundary compliance checks.

\item[Failure Mode and Effects Analysis (FMEA).] A structured, systematic methodology for enumerating the potential failure modes of a system, assessing their root causes and downstream consequences, and quantifying associated risk~\cite{iec60812}. This framework extends the conventional FMEA to encompass probabilistic and adversarial failure modes specific to AI systems, referred to herein as a \textit{model-focused FMEA}.

\end{deflist}
\acused{AI}\acused{DCE}\acused{OOD}\acused{HITL}\acused{CICD}\acused{FMEA}

\needspace{25\baselineskip}
\begin{samepage}
\section{The Five Core Engineering Principles}
\label{sec:principles}

\subsection*{Principle I - Verification}
You shall not accept any algorithmic predictions or outputs as being verified truths. All data models, sensitivities, optimisations, and simulations must be constrained within known engineering, statistical, and physical limits, and rigorously validated against empirical reality before use.


\subsection*{Principle II - Traceability}
Provenance of training datasets and generated data must be evident to ensure the pipeline is fully reproducible and auditable. For third-party or open-source models where raw provenance is inaccessible, a formal secondary data attestation documenting known dataset limitations is required. Misrepresenting the source of data, concealing algorithmic bias, or promoting unverified AI-generated content as your own engineering design is unethical.


\subsection*{Principle III - Ultimate Liability}
The algorithm does not possess professional liability; only the engineer does. You have ultimate, nondelegable liability for the physical and clinical failure modes of any system powered by an AI pipeline if you certify it.

\subsection*{Principle IV - Boundary Clarity}
It is your responsibility to define explicitly where the training distribution ends and extrapolation begins using mathematically defensible metrics (e.g., conformal prediction intervals or reconstruction error thresholds). You must foster an environment in which model anomalies, distributional drift, and \ac{OOD} conditions are systematically detected and disclosed.

\subsection*{Principle V - Stewardship}
In fields where technology is advancing rapidly, you must anticipate future failure modes, adversarial vulnerabilities, and the potential for misuse or dual-use of your systems. Where there exists a credible risk of serious harm to human beings or the natural environment, the precautionary principle shall govern.

\end{samepage}


\section{Operationalizing the Framework}
\label{sec:operationalizing}

To avoid the danger of keeping this framework purely philosophical in nature, it is imperative that these principles be operationalized into concrete engineering processes. This involves incorporating verification checkpoints into the technical system.

\subsection{Proportional Governance Using Tiers of Risk}
To ensure institutional feasibility, compliance constraints are applied proportionally via a three-tiered risk framework rather than a blanket mandate.

\begin{table}[H]
\centering
\caption{Data-Centric Engineering Risk Classification and Governance Matrix}
\label{tab:risk_tiers}
\vspace{2mm}
\small
\begin{tabularx}{\textwidth}{l X X}
\toprule
\textbf{Risk Tier} & \textbf{Verification Boundary \& Operational Profile} & \textbf{Mandatory Governance \& Validation Requirements} \\
\midrule
\textbf{Tier 1: Low} & Computational pathways are entirely deterministic and transparent; the practitioner possesses the analytical capacity to independently audit, reproduce, and validate the specific mechanisms generating the output. & Exempt from formal external oversight. The practitioner is required to maintain comprehensive documentation detailing foundational assumptions, data lineage, and validation methodologies within established institutional records. \\
\midrule
\textbf{Tier 2: Medium} & System complexity or dimensionality prevents direct, step-by-step cognitive auditing of individual internal operations, requiring the practitioner to evaluate system performance through aggregate behavioral characteristics and bounded validation metrics. & Mandatory quantification of operational margins and uncertainty; formalized procedures for expert intervention and qualitative appraisal; documented internal technical verification prior to downstream utilization. \\
\midrule
\textbf{Tier 3: High} & System execution exhibits complete operational opacity or inherent non-determinism, preventing the practitioner from predicting or validating the accuracy of specific outputs prior to operational deployment and systemic integration. & Implementation of independent, structurally isolated boundary constraints capable of overriding system execution; exhaustive evaluation under simulated boundary conditions; formal acknowledgement explicitly accepting the epistemic risk of deploying unverified computational outputs. \\
\bottomrule
\end{tabularx}
\end{table}

\subsection{Integration into the Engineering Lifecycle}
This framework establishes a multi-layered governance model that departs from conventional, retrospective auditing paradigms. Rather than serving as a post-hoc diagnostic tool, compliance is operationalized as an active, continuous gating mechanism integrated throughout the engineering lifecycle across three sequential phases.

\paragraph{Project Inception.}
Governance controls are mandated prior to architectural design. Each computational pipeline must clear an established project inception gate requiring a rigorous structural risk assessment. Executing this evaluation at the pre-development stage explicitly delineates the operational scope of the system, maps potential failure modes, and establishes formal mitigation protocols before resources are committed.

\paragraph{Continuous Automated Pipeline Verification.}
During development, governance requirements are translated into automated verification protocols embedded within the compilation and integration pipeline. Alongside conventional deterministic unit testing, the pipeline executes automated edge-case evaluations, data lineage verification, and mathematical boundary checks. Should an iteration violate established physical laws or exhibit statistical drift from its baseline distribution, compilation is automatically halted, preventing deployment.

\paragraph{Decoupled System Architecture.}
The deployment phase mandates a structurally decoupled architecture for runtime execution within high-stakes environments. The probabilistic computational pipeline operates strictly as an advisory system, subordinated to a deterministic, non-probabilistic safety loop. Upon detection of \ac{OOD} data inputs or operational anomalies, the architecture triggers an immediate isolation protocol. This structural separation guarantees the instantaneous fallback of the physical asset into a predetermined, deterministically safe state, requiring formal expert validation and \ac{HITL} override before autonomous operations can resume.

\subsection{Academic Culture and Research Training}
Technical countermeasures remain insufficient without an underlying institutional culture of systematic skepticism. Academic curricula and research training protocols must condition practitioners to treat computational outputs not as objective verities, but as complex conditional probability distributions requiring rigorous empirical validation.

\subsection{Practical Guidelines for Effective AI Use}
AI is appropriate when the output is verifiable, AI must be seen as a tool and when use correctly it can accelerate workflow and handle tasks however it must never be seen as a substitution for a human or a human's judgement. Efficient and Effective Use (EEU) of AI is possible and encouraged both in industry and academia and stems from being able to determine where and when to use it and understanding the cost of using it. Costs of AI can lie within the time to evaluate its output, the ethics of misattribution of computationally generated material as human scholarship and the epistemic risk of allowing AI to make inferences that require a humans judgement. Both students and Faculty are capable of EEU and should be encouraged in these safe practices but they should remain eternally aware of the risks of stochastic artifacts, bias and unquantified uncertainty present in AI systems.

\section{Conclusion}
\label{sec:conclusion}
% Reiterate that accountability cannot be outsourced to code. By enforcing strict bounding constraints, explicit data provenance, clear boundary controls, and unyielding human professional liability, the organization can safely deploy cutting-edge AI pipelines without compromising its fundamental dedication to public safety and physical engineering excellence.
Accountability cannot be outsourced. Automated coding and AI-assisted tooling can accelerate development, but no tool can substitute the professional responsibility of ensuring a system is safe, validated, and fit for deployment in safety-critical environments. It is the duty of an engineer to enforce strict boundary constraints, explicit data provenance, clear boundary controls, and to acknoledge the unyielding human professional liability. By doing so, this encourages the safe and effective usage of AI without compromising the fundamental dedication to public safety and physical engineering excellence.

\section*{List of Acronyms}
\begin{acronym}[HITL]
  \acro{AI}{Artificial Intelligence}
  \acro{ASCE}{American Society of Civil Engineers}
  \acro{ASME}{American Society of Mechanical Engineers}
  \acro{CICD}[CI/CD]{Continuous Integration/Continuous Deployment}
  \acro{DCE}{Data-Centric Engineering}
  \acro{FMEA}{Failure Mode and Effects Analysis}
  \acro{HITL}{Human-in-the-Loop}
  \acro{IEC}{International Electrotechnical Commission}
  \acro{IEEE}{Institute of Electrical and Electronics Engineers}
  \acro{ISO}{International Organization for Standardization}
  \acro{OOD}{Out-of-Distribution}
\end{acronym}

\printbibliography

\end{document}
